\section{INTRODUÇÃO}

A presença de animais de estimação tornou-se cada vez mais comum nos lares brasileiros, consolidando os animais de companhia como verdadeiros membros do núcleo familiar. Essa estreita relação de afeto faz com que a convivência diária transforme a rotina das pessoas, mas também traz preocupações constantes com o bem-estar e a segurança dos animais. Nesse contexto, episódios inesperados de fuga e desaparecimento causam intenso abalo emocional aos tutores e exigem ações de busca imediatas e articuladas na vizinhança.

Contudo, a perda de animais nas cidades configura um problema urbano complexo e de difícil enfrentamento. Quando um animal escapa para as vias públicas, a ausência de um sistema centralizado e georreferenciado dificulta a identificação rápida do seu paradeiro, elevando os riscos de acidentes e sobrecarregando protetores e abrigos locais. Historicamente, a comunicação de desaparecimentos dependia de cartazes de papel colados em postes e avisos entre vizinhos, práticas que com o tempo migraram para postagens espalhadas em redes sociais. Embora alcancem grande audiência, esses canais não foram projetados para gestão de emergências, pois as publicações perdem visibilidade rapidamente nas linhas do tempo e não dispõem de visualização sobre mapas cartográficos. No cenário atual, há uma grande falta de ferramentas gratuitas, colaborativas e focadas na localização espacial imediata do animal.

Diante desse cenário, este trabalho apresenta o desenvolvimento do sistema WeFIND, um aplicativo móvel colaborativo criado para centralizar e agilizar a localização e divulgação de animais perdidos e encontrados, além de apoiar a adoção responsável. A plataforma organiza as publicações sobre um mapa interativo em tempo real, permitindo identificar visualmente animais na vizinhança. Para qualificar as buscas, o sistema padroniza o cadastro com fotos e traços essenciais do animal como espécie, raça, porte e pelagem, oferece um canal de mensagens privativo entre os usuários, automatiza a geração de cartazes com código QR para impressão física e disponibiliza o cadastro preventivo de animais tutelados com código QR para identificação imediata.

O objetivo central desta pesquisa é desenvolver uma aplicação móvel colaborativa que auxilie tutores e a comunidade na localização e divulgação de animais perdidos e encontrados. Para alcançar esse propósito, o trabalho estrutura-se no cumprimento de metas específicas que envolvem a realização de uma pesquisa prática com a comunidade para identificar as principais dificuldades na busca por animais, a avaliação técnica de sistemas similares existentes no mercado, o levantamento e a documentação detalhada dos requisitos funcionais e não funcionais, a elaboração da modelagem conceitual em notação UML, o projeto da interface centrado na experiência do usuário, a implementação do sistema utilizando React Native e Supabase e, por fim, a realização de testes funcionais e de usabilidade com voluntários do público-alvo.

A relevância social deste projeto está em fornecer uma ferramenta prática e acessível para diminuir o tempo de busca e apoiar a causa animal. Do ponto de vista técnico e acadêmico, o trabalho justifica-se pela oportunidade de integrar conhecimentos de Engenharia de Software, programação para dispositivos móveis, bancos de dados e design de interface desenvolvidos ao longo do Curso Superior de Tecnologia em Sistemas para Internet do IFSul.

A metodologia adotada combina pesquisa descritiva com desenvolvimento tecnológico baseado no modelo incremental, utilizando o método ágil Kanban por meio da ferramenta Trello para o acompanhamento sistemático de cada etapa.

O restante desta monografia está organizado em capítulos sequenciais. O Capítulo 2 contextualiza o crescimento da população de animais domésticos no Brasil, as problemáticas decorrentes de fugas e as limitações de métodos como colar cartazes de papel em postes e publicar em grupos de redes sociais. O Capítulo 3 detalha os procedimentos metodológicos da pesquisa de campo com tutores, o estudo comparativo de sistemas similares e a metodologia ágil de desenvolvimento de software. O Capítulo 4 formaliza a engenharia de requisitos funcionais e não funcionais e a modelagem conceitual do sistema WeFIND em diagramas da UML. O Capítulo 5 apresenta a concepção da plataforma WeFIND, a motivação empírica pessoal que originou o projeto, seus quatro pilares conceituais principais e a justificativa da abordagem móvel em relação à web. O Capítulo 6 descreve as tecnologias do ecossistema de desenvolvimento e a modelagem lógica relacional do banco de dados PostgreSQL com políticas de segurança RLS no Supabase. O Capítulo 7 aborda o projeto de design de interface centrado no usuário, detalhando a tipografia, a identidade visual e as telas do aplicativo com análise detalhada de seus elementos. O Capítulo 8 documenta o planejamento, a execução e os resultados dos testes funcionais de caixa-preta e da avaliação de usabilidade pela escala SUS. Por fim, o Capítulo 9 reúne as considerações finais, as dificuldades técnicas superadas e as propostas de trabalhos futuros.
\clearpage
